% Part: sets-functions-relations
% Chapter: relations
% Section: reflections
%
\documentclass[../../../include/open-logic-section]{subfiles}

\begin{document}

\olfileid[fa-IR]{sfr}{rel}{ref}
\olsection{تأملات فلسفی}

در \olref[set]{sec}، روابط را مجموعه‌هایی معین تعریف کردیم. شایسته است
درنگ کنیم و پرسشی فلسفی و کوتاه بپرسیم: چنین تعریفی چه
\emph{می‌کند}؟ بسیار بعید است بخواهیم بگوییم برخی واقعیت‌های متافیزیکیِ
این‌همانی را \emph{کشف کرده‌ایم}؛ مثلاً اینکه رابطهٔ ترتیب روی $\Nat$،
چنان‌که \emph{معلوم شده است}، همان مجموعهٔ
$R=  \Setabs{\tuple{n,m}}{n, m \in \Nat\text{ و }
n < m}$ است که در \olref[set]{sec} تعریف کردیم. در ادامه سه دلیل
برای این تردید می‌آوریم.

نخست: در \olref[set][pai]{wienerkuratowski}، تعریف کردیم $\tuple{a, b} =
\{\{a\}, \{a, b\}\}$. در عوض، تعریف $\lVert a,
b\rVert = \{\{b\}, \{a, b\}\} = \tuple{b,a}$ را در نظر بگیرید. وقتی
$a \neq b$، داریم $\tuple{a, b} \neq \lVert a,b\rVert$. اما به همان
اندازه می‌توانستیم $\lVert a,b\rVert$ را تعریف خود از زوج مرتب قرار
دهیم، نه $\tuple{a,b}$ را. هر دو تعریف به یک اندازه کار می‌کردند. پس
اکنون دو نامزدِ به یک اندازه شایسته داریم که هر یک می‌تواند «خودِ»
رابطهٔ ترتیب روی اعداد طبیعی باشد، یعنی:
\begin{align*}
		R &= \Setabs{\tuple{n,m}}{n, m \in \Nat \text{ و }n < m}\\
		S &= \Setabs{\lVert n,m\rVert}{n, m \in \Nat \text{ و }n < m}.
\end{align*}
از آنجا که بنا بر اصل گسترش $R \neq S$، روشن است که نمی‌توانند
\emph{هر دو} با رابطهٔ ترتیب روی~$\Nat$ این‌همان باشند. اما ادعای اینکه
بنا بر واقعیت امر، $R$، و نه $S$ (یا برعکس)، \emph{همان} رابطهٔ ترتیب
است، صرفاً دل‌بخواهی و ازاین‌رو اندکی شرم‌آور خواهد بود. (این نمونه‌ای بسیار ساده
از استدلالی علیه تقلیل‌گرایی مجموعه‌شناختی است که بناسراف در
\citeyear{Benacerraf1965} مشهور ساخت. چند بار دیگر به آن بازخواهیم گشت.)

دوم: اگر گمان کنیم \emph{هر} رابطه‌ای باید با مجموعه‌ای یکی انگاشته
شود، آنگاه خودِ رابطهٔ عضویت در مجموعه، یعنی $\in$، نیز باید مجموعه‌ای
معین باشد. در واقع، باید مجموعهٔ $\Setabs{\tuple{x,y}}{x \in y}$ باشد.
اما آیا این مجموعه وجود دارد؟ با توجه به پارادوکس راسل، وجود چنین
مجموعه‌ای ادعایی نابدیهی است. در حقیقت،
\oliflabeldef{cumul:::part}{نظریهٔ مجموعه‌ها که آن را در
\olref[cumul][][]{part} بسط می‌دهیم، وجود این مجموعه را \emph{انکار}
خواهد کرد.\footnote{اگر قدری جلوتر برویم، دلیل چنین است. برای برهان
خلف، فرض کنید $I =
\Setabs{\tuple{x,y}}{x \in y}$ وجود دارد. آنگاه $\bigcup \bigcup I$
مجموعهٔ جهان‌شمول است، و این با
\olref[sfr][z][sep]{thm:NoUniversalSet} تناقض دارد.}}{می‌توان نظریهٔ
مجموعه‌ها را به شیوه‌ای دقیق و در قالب نظریه‌ای اصل‌موضوعی بسط داد، و
آن نظریه واقعاً وجود این مجموعه را انکار خواهد کرد.}
پس حتی اگر بتوان برخی روابط را همچون مجموعه‌ها در نظر گرفت، رابطهٔ
عضویت در مجموعه باید حالتی ویژه باشد.

سوم: وقتی روابط را با مجموعه‌ها «یکی می‌انگاریم»، گفتیم به خود اجازه
می‌دهیم $Rxy$ را به جای $\tuple{x,y} \in R$ بنویسیم. این کار بی‌اشکال
است، به شرط آنکه رابطهٔ عضویت، یعنی «$\in$»، \emph{به‌منزلهٔ} یک محمول
در نظر گرفته شود. اما اگر گمان کنیم «$\in$» بر نوع معینی از مجموعه
دلالت دارد، آنگاه عبارت «$\tuple{x,y} \in R$» فقط از سه حد مفرد تشکیل
شده است که بر مجموعه‌ها دلالت می‌کنند: «$\tuple{x,y}$»، «$\in$» و
«$R$». چنین فهرستی از نام‌ها، درست مانند رشتهٔ بی‌معنای «فنجان جاقلمی
میز»، توان بیان یک گزاره را ندارد. باز هم، حتی اگر بتوان برخی روابط را همچون
مجموعه‌ها در نظر گرفت، رابطهٔ عضویت در مجموعه باید حالتی ویژه باشد.
(این استدلال صورت ساده‌ای از پارادوکس «مفهوم \emph{اسب}» فرگه را با
اعتراض مشهوری درهم می‌آمیزد که ویتگنشتاین زمانی علیه راسل مطرح کرد.)

پس این بحث ما را به کجا می‌رساند؟ خب، اینکه بگوییم روابط روی اعداد
مجموعه‌اند هیچ \emph{اشکالی} ندارد. فقط باید روح حاکم بر این سخن را
دریابیم. ما واقعیتی متافیزیکی دربارهٔ این‌همانی بیان نمی‌کنیم؛ صرفاً
خاطرنشان می‌کنیم که در برخی زمینه‌ها می‌توانیم (و چنین خواهیم کرد)
روابطی (معین) را \emph{همچون} مجموعه‌هایی معین در نظر بگیریم.

\end{document}
